% ============================================================
% Question Bank: ch08-01 Populations and Samples
% Topic Title: Populations and Samples
% CCSS: 7.SP.A.1
% Questions: 30 (18 MC + 7 SA + 5 Graphical)
% ============================================================

% --- Multiple Choice Questions (18) ---

\begin{practiceQuestion}{8-1-q01}{mc}
\begin{questionText}
A school has $800$ students. The principal surveys $50$ students about their favorite subject. What is the population?
\end{questionText}
\choiceA{The $50$ surveyed students}
\choiceB{All $800$ students in the school}
\choiceC{The principal}
\choiceD{The students who like math}
\correctAnswer{B}
\explanation{The population is the entire group you want to study — all $800$ students.}
\end{practiceQuestion}

\begin{practiceQuestion}{8-1-q02}{mc}
\begin{questionText}
A researcher wants to know the average height of seventh graders in a state. She measures the heights of $100$ seventh graders from different schools. What is the sample?
\end{questionText}
\choiceA{All seventh graders in the state}
\choiceB{All students in the schools she visited}
\choiceC{The $100$ seventh graders she measured}
\choiceD{The tallest seventh graders}
\correctAnswer{C}
\explanation{The sample is the smaller group chosen from the population — the $100$ students measured.}
\end{practiceQuestion}

\begin{practiceQuestion}{8-1-q03}{mc}
\begin{questionText}
Which sampling method is most likely to produce a representative sample of a school's students?
\end{questionText}
\choiceA{Survey students in the cafeteria during lunch}
\choiceB{Survey every fifth student on an alphabetical class list}
\choiceC{Survey students who volunteer to participate}
\choiceD{Survey only students in honors classes}
\correctAnswer{B}
\explanation{Selecting every fifth student from a list is systematic and gives all students an equal chance, making it more representative.}
\end{practiceQuestion}

\begin{practiceQuestion}{8-1-q04}{mc}
\begin{questionText}
A company wants to know if customers like a new product. They survey only customers who bought the product last week. What is the problem with this sample?
\end{questionText}
\choiceA{The sample is too large}
\choiceB{The sample is biased toward recent buyers}
\choiceC{The sample is random}
\choiceD{There is no problem}
\correctAnswer{B}
\explanation{Only surveying recent buyers introduces bias — they may feel differently than all customers.}
\end{practiceQuestion}

\begin{practiceQuestion}{8-1-q05}{mc}
\begin{questionText}
A random sample means that:
\end{questionText}
\choiceA{Only the best members are chosen}
\choiceB{The largest possible group is chosen}
\choiceC{Every member of the population has an equal chance of being selected}
\choiceD{Only volunteers are selected}
\correctAnswer{C}
\explanation{In a random sample, every member of the population has an equal chance of being chosen.}
\end{practiceQuestion}

\begin{practiceQuestion}{8-1-q06}{mc}
\begin{questionText}
A city wants to survey residents about a new park. Which method would produce the least biased sample?
\end{questionText}
\choiceA{Survey people at the existing park}
\choiceB{Put the survey in the newspaper and wait for responses}
\choiceC{Randomly select addresses from the city directory}
\choiceD{Survey people at the city council meeting}
\correctAnswer{C}
\explanation{Randomly selecting addresses gives all residents an equal chance, reducing bias.}
\end{practiceQuestion}

\begin{practiceQuestion}{8-1-q07}{mc}
\begin{questionText}
Which of the following is an example of a biased sample?
\end{questionText}
\choiceA{Randomly selecting $50$ names from a hat containing all student names}
\choiceB{Using a random number generator to pick $30$ employees}
\choiceC{Asking every $10$th person who enters a store}
\choiceD{Surveying only members of the basketball team about school sports}
\correctAnswer{D}
\explanation{Basketball players have strong opinions about sports that may not represent the whole school.}
\end{practiceQuestion}

\begin{practiceQuestion}{8-1-q08}{mc}
\begin{questionText}
A sample of $40$ out of $2{,}000$ voters is surveyed. What fraction of the population is in the sample?
\end{questionText}
\choiceA{$\frac{1}{50}$}
\choiceB{$\frac{1}{40}$}
\choiceC{$\frac{1}{20}$}
\choiceD{$\frac{1}{500}$}
\correctAnswer{A}
\explanation{$\frac{40}{2{,}000} = \frac{1}{50}$.}
\end{practiceQuestion}

\begin{practiceQuestion}{8-1-q09}{mc}
\begin{questionText}
A teacher wants to know how many hours per week students in her school read. She surveys all students in her own class. Why might this not represent the whole school?
\end{questionText}
\choiceA{Her class is too large}
\choiceB{Her class may not be typical of the whole school}
\choiceC{Students in her class read more than average}
\choiceD{She should survey fewer students}
\correctAnswer{B}
\explanation{One class may have different reading habits than other classes, so it may not represent the whole school.}
\end{practiceQuestion}

\begin{practiceQuestion}{8-1-q10}{mc}
\begin{questionText}
Which question would require sampling rather than surveying the entire population?
\end{questionText}
\choiceA{How many students are in a classroom of $25$?}
\choiceB{What is the average lifespan of a light bulb produced by a factory?}
\choiceC{What is the total number of books on a shelf?}
\choiceD{How many chairs are in a room?}
\correctAnswer{B}
\explanation{Testing every light bulb would destroy them all. Sampling lets you estimate the average without testing every one.}
\end{practiceQuestion}

\begin{practiceQuestion}{8-1-q11}{mc}
\begin{questionText}
A survey of $200$ randomly selected adults found that $60\%$ prefer coffee over tea. If the city has $50{,}000$ adults, about how many prefer coffee?
\end{questionText}
\choiceA{$12{,}000$}
\choiceB{$20{,}000$}
\choiceC{$30{,}000$}
\choiceD{$60{,}000$}
\correctAnswer{C}
\explanation{$60\%$ of $50{,}000 = 0.60 \times 50{,}000 = 30{,}000$.}
\end{practiceQuestion}

\begin{practiceQuestion}{8-1-q12}{mc}
\begin{questionText}
Which of the following increases the reliability of conclusions drawn from a sample?
\end{questionText}
\choiceA{Using a smaller sample}
\choiceB{Using a larger random sample}
\choiceC{Surveying only people who agree with you}
\choiceD{Using a non-random sample}
\correctAnswer{B}
\explanation{Larger random samples better represent the population and give more reliable results.}
\end{practiceQuestion}

\begin{practiceQuestion}{8-1-q13}{mc}
\begin{questionText}
A manager surveys every $20$th customer who enters a store. This is an example of:
\end{questionText}
\choiceA{A voluntary sample}
\choiceB{A convenience sample}
\choiceC{A systematic sample}
\choiceD{A biased sample}
\correctAnswer{C}
\explanation{Selecting every $n$th member is called systematic sampling.}
\end{practiceQuestion}

\begin{practiceQuestion}{8-1-q14}{mc}
\begin{questionText}
A website asks visitors to rate a product. Only $5\%$ of visitors respond. This is a:
\end{questionText}
\choiceA{Random sample}
\choiceB{Systematic sample}
\choiceC{Voluntary response sample}
\choiceD{Stratified sample}
\correctAnswer{C}
\explanation{When people choose whether to participate, it is a voluntary response sample, which tends to be biased.}
\end{practiceQuestion}

\begin{practiceQuestion}{8-1-q15}{mc}
\begin{questionText}
In a school of $600$ students, a random sample of $60$ found that $15$ students walk to school. Predict how many students in the entire school walk to school.
\end{questionText}
\choiceA{$100$}
\choiceB{$150$}
\choiceC{$200$}
\choiceD{$250$}
\correctAnswer{B}
\explanation{$\frac{15}{60} = \frac{1}{4} = 25\%$. So $25\%$ of $600 = 150$.}
\end{practiceQuestion}

\begin{practiceQuestion}{8-1-q16}{mc}
\begin{questionText}
Why is a convenience sample (e.g., surveying your friends) usually not a good way to learn about a population?
\end{questionText}
\choiceA{It is too expensive}
\choiceB{It takes too long}
\choiceC{It is likely biased and not representative}
\choiceD{It includes too many people}
\correctAnswer{C}
\explanation{Convenience samples are easy but likely biased because they don't give everyone an equal chance.}
\end{practiceQuestion}

\begin{practiceQuestion}{8-1-q17}{mc}
\begin{questionText}
A doctor wants to study the effect of a new medicine on $10{,}000$ patients. She randomly selects $500$ patients for a trial. What is the sample size?
\end{questionText}
\choiceA{$10{,}000$}
\choiceB{$500$}
\choiceC{$10{,}500$}
\choiceD{$50$}
\correctAnswer{B}
\explanation{The sample size is the number of individuals selected — $500$ patients.}
\end{practiceQuestion}

\begin{practiceQuestion}{8-1-q18}{mc}
\begin{questionText}
A poll found that $35\%$ of a random sample of $400$ students prefer pizza for lunch. Which statement is the best inference?
\end{questionText}
\choiceA{Exactly $35\%$ of all students prefer pizza}
\choiceB{About $35\%$ of the entire student population likely prefers pizza}
\choiceC{No students prefer anything other than pizza}
\choiceD{The poll is invalid because not all students were surveyed}
\correctAnswer{B}
\explanation{A random sample provides an estimate. About $35\%$ of the population likely prefers pizza, but the exact percentage may differ slightly.}
\end{practiceQuestion}

% --- Short Answer Questions (7) ---

\begin{practiceQuestion}{8-1-q19}{sa}
\begin{questionText}
A random sample of $80$ out of $2{,}000$ students found that $20$ students ride the bus. Predict how many students in the school ride the bus.
\end{questionText}
\correctAnswer{$500$}
\explanation{$\frac{20}{80} = 25\%$. Then $25\%$ of $2{,}000 = 500$.}
\end{practiceQuestion}

\begin{practiceQuestion}{8-1-q20}{sa}
\begin{questionText}
Define the term \emph{population} in statistics.
\end{questionText}
\correctAnswer{The entire group you want to study}
\explanation{The population is the whole group about which you want to draw conclusions.}
\end{practiceQuestion}

\begin{practiceQuestion}{8-1-q21}{sa}
\begin{questionText}
A random sample of $50$ apples from an orchard found that $6$ were bruised. If the orchard has $3{,}000$ apples, about how many are bruised?
\end{questionText}
\correctAnswer{$360$}
\explanation{$\frac{6}{50} = 12\%$. Then $12\%$ of $3{,}000 = 360$.}
\end{practiceQuestion}

\begin{practiceQuestion}{8-1-q22}{sa}
\begin{questionText}
Give one reason why a sample might not perfectly represent a population.
\end{questionText}
\correctAnswer{Random variation (or the sample may be biased)}
\explanation{Even a well-chosen random sample may differ slightly from the population due to chance.}
\end{practiceQuestion}

\begin{practiceQuestion}{8-1-q23}{sa}
\begin{questionText}
A radio station asks listeners to call in and vote for their favorite song. Explain why this is not a random sample.
\end{questionText}
\correctAnswer{Only listeners who choose to call participate, so it is voluntary and biased}
\explanation{Voluntary response samples are biased because people with strong opinions are more likely to respond.}
\end{practiceQuestion}

\begin{practiceQuestion}{8-1-q24}{sa}
\begin{questionText}
A random sample of $100$ out of $4{,}000$ residents found that $45$ support a new library. Predict how many residents support it.
\end{questionText}
\correctAnswer{$1{,}800$}
\explanation{$\frac{45}{100} = 45\%$. Then $45\%$ of $4{,}000 = 1{,}800$.}
\end{practiceQuestion}

\begin{practiceQuestion}{8-1-q25}{sa}
\begin{questionText}
A factory produces $10{,}000$ batteries per day. A sample of $200$ batteries is tested, and $8$ are defective. What percent of the sample is defective?
\end{questionText}
\correctAnswer{$4\%$}
\explanation{$\frac{8}{200} = 0.04 = 4\%$.}
\end{practiceQuestion}

% --- Graphical Questions (3 GMC + 2 GSA) ---

\begin{practiceQuestion}{8-1-q26}{gmc}
\begin{questionText}
A student surveyed $50$ classmates about their favorite fruit. The results are shown in the bar graph below.

\begin{center}
\begin{tikzpicture}[scale=0.7]
  \draw[thick,->] (0,0) -- (0,6.5) node[above]{\small Number of Students};
  \draw[thick,->] (0,0) -- (8.5,0);
  % bars
  \fill[funBlue!70] (0.5,0) rectangle (1.5,4);
  \fill[funGreen!70] (2.5,0) rectangle (3.5,3);
  \fill[funOrange!70] (4.5,0) rectangle (5.5,2);
  \fill[funRed!70] (6.5,0) rectangle (7.5,1);
  % labels
  \node[below] at (1,0) {\small Apple};
  \node[below] at (3,0) {\small Banana};
  \node[below] at (5,0) {\small Orange};
  \node[below] at (7,0) {\small Grape};
  % y-axis ticks
  \foreach \y in {1,2,3,4,5,6} {
    \draw (-0.1,\y) -- (0.1,\y);
    \node[left] at (0,\y) {\small $\y\times5$};
  }
  % values on bars
  \node[above] at (1,4) {\small $20$};
  \node[above] at (3,3) {\small $15$};
  \node[above] at (5,2) {\small $10$};
  \node[above] at (7,1) {\small $5$};
\end{tikzpicture}
\end{center}

If there are $500$ students in the school, predict how many prefer apples.
\end{questionText}
\choiceA{$20$}
\choiceB{$100$}
\choiceC{$200$}
\choiceD{$250$}
\correctAnswer{C}
\explanation{$\frac{20}{50} = 40\%$ prefer apples. $40\%$ of $500 = 200$.}
\end{practiceQuestion}

\begin{practiceQuestion}{8-1-q27}{gmc}
\begin{questionText}
The diagram below shows four different sampling methods used to survey students at a school. Which method is most likely to produce an unbiased sample?

\begin{center}
\begin{tikzpicture}[scale=0.75]
  % Method A
  \node[draw,rounded corners,fill=funBlue!15,minimum width=3.5cm,minimum height=1.2cm,align=center] (A) at (0,3) {\textbf{Method A}\\Survey the football team};
  % Method B
  \node[draw,rounded corners,fill=funGreen!15,minimum width=3.5cm,minimum height=1.2cm,align=center] (B) at (6,3) {\textbf{Method B}\\Survey every 10th student\\on the school roster};
  % Method C
  \node[draw,rounded corners,fill=funOrange!15,minimum width=3.5cm,minimum height=1.2cm,align=center] (C) at (0,0) {\textbf{Method C}\\Post survey online and\\wait for responses};
  % Method D
  \node[draw,rounded corners,fill=funRed!15,minimum width=3.5cm,minimum height=1.2cm,align=center] (D) at (6,0) {\textbf{Method D}\\Survey students in\\one math class};
\end{tikzpicture}
\end{center}
\end{questionText}
\choiceA{Method A}
\choiceB{Method B}
\choiceC{Method C}
\choiceD{Method D}
\correctAnswer{B}
\explanation{Selecting every 10th student from the school roster is systematic and gives all students a chance, making it the least biased.}
\end{practiceQuestion}

\begin{practiceQuestion}{8-1-q28}{gmc}
\begin{questionText}
The table below shows the results of two samples taken from a population of $1{,}000$ students.

\begin{center}
\begin{tabular}{|l|c|c|}
\hline
 & \textbf{Sample 1 (50 students)} & \textbf{Sample 2 (50 students)} \\
\hline
Prefer Math & $18$ & $22$ \\
\hline
Prefer Science & $14$ & $12$ \\
\hline
Prefer English & $10$ & $8$ \\
\hline
Prefer Art & $8$ & $8$ \\
\hline
\end{tabular}
\end{center}

Based on both samples, the best prediction for how many of the $1{,}000$ students prefer math is:
\end{questionText}
\choiceA{$180$}
\choiceB{$220$}
\choiceC{$400$}
\choiceD{$360$}
\correctAnswer{C}
\explanation{Average the two samples: $\frac{18+22}{2} = 20$ out of $50 = 40\%$. Then $40\%$ of $1{,}000 = 400$.}
\end{practiceQuestion}

\begin{practiceQuestion}{8-1-q29}{gsa}
\begin{questionText}
The bar graph shows the results of surveying a random sample of $60$ students about how they get to school.

\begin{center}
\begin{tikzpicture}[scale=0.65]
  \draw[thick,->] (0,0) -- (0,6) node[above]{\small Students};
  \draw[thick,->] (0,0) -- (7,0);
  \fill[funBlue!70] (0.5,0) rectangle (1.5,5);
  \fill[funGreen!70] (2.5,0) rectangle (3.5,3);
  \fill[funOrange!70] (4.5,0) rectangle (5.5,2);
  \node[below,font=\small] at (1,-0.1) {Bus};
  \node[below,font=\small] at (3,-0.1) {Walk};
  \node[below,font=\small] at (5,-0.1) {Car};
  \foreach \y in {1,...,5} {
    \draw (-0.1,\y) -- (0.1,\y);
    \node[left,font=\small] at (0,\y) {$\y\times6$};
  }
  \node[above,font=\small] at (1,5) {$30$};
  \node[above,font=\small] at (3,3) {$18$};
  \node[above,font=\small] at (5,2) {$12$};
\end{tikzpicture}
\end{center}

If the school has $900$ students, predict how many walk to school.
\end{questionText}
\correctAnswer{$270$}
\explanation{$\frac{18}{60} = 30\%$. Then $30\%$ of $900 = 270$.}
\end{practiceQuestion}

\begin{practiceQuestion}{8-1-q30}{gsa}
\begin{questionText}
The table shows the results of a random sample of $80$ people about their preferred pet.

\begin{center}
\begin{tabular}{|l|c|}
\hline
\textbf{Pet} & \textbf{Number} \\
\hline
Dog & $32$ \\
\hline
Cat & $24$ \\
\hline
Fish & $12$ \\
\hline
Bird & $8$ \\
\hline
Other & $4$ \\
\hline
\end{tabular}
\end{center}

What percent of the sample prefers cats? If the population is $5{,}000$, predict how many prefer cats.
\end{questionText}
\correctAnswer{$30\%$; $1{,}500$ people}
\explanation{$\frac{24}{80} = 0.30 = 30\%$. Then $30\%$ of $5{,}000 = 1{,}500$.}
\end{practiceQuestion}
